% 第16批：原书印刷第116—117页；源图题标注复核条件。
\subsection{自测题5：理想气体和热力学填空}

\sourcetest{5.2.1}{116}{124}
\problemtext 氢分子质量 $3.3\times10^{-24}$g，每秒有 $10^{23}$ 个分子以 $10^5$cm/s 速度、与壁面法线夹45°，在 $2.0$cm$^2$ 面积内完全弹性反弹；求压强。
\workings 单次法向冲量 $2mv\cos45^\circ$，$p=N(2mv\cos45^\circ)/(A\Delta t)$。代入 $m=3.3\times10^{-27}$kg、$v=10^3$m/s、$A=2.0\times10^{-4}$m$^2$、$N/\Delta t=10^{23}$s$^{-1}$，得 $\boxed{p\simeq2.33\times10^3\ {\rm Pa}}$。
\crosscheck 法线方向速度是 $v\cos45^\circ$；不能把单次碰壁动量变化写成 $mv$。

\sourcetest{5.2.2}{116}{124}
\problemtext 理想气体密度在等压升温、等温增压两种过程中分别怎样改变？
\workings 由 $\rho=pM/(RT)$，等压时 $\rho\propto1/T$，\emph{减小}；等温时 $\rho\propto p$，\emph{增大}。
\crosscheck 摄氏温度不可直接代入理想气体绝对温度比例。

\sourcetest{5.2.3}{116}{124}
\problemtext 氢气作为刚性双原子理想气体，升温1K时内能增加41.6J；求气体质量。
\workings $\Delta U=n(5R/2)\Delta T$，故 $n=41.6/(2.5\times8.31)\simeq2.00$mol。取氢气摩尔质量2g/mol，$\boxed{m\simeq4.00\ {\rm g}}$。
\crosscheck 采用常温转动自由度已激发、振动自由度未激发的刚性双原子近似。

\sourcetest{5.2.4}{117}{125}
\problemtext 定量理想气体先等容加热至原来绝对温度4倍，后等温膨胀至原来体积2倍；求分子平均撞击器壁频率倍率。
\workings 撞壁频率尺度 $z\propto n\bar v\propto (N/V)\sqrt T$，终态体积加倍、温度变4倍，故 $z_2/z_1=(1/2)\cdot2=\boxed{1}$。
\crosscheck 在固定截面积或单位面积的撞壁频率意义下成立；并非单个分子对随动器壁的碰撞计数。

\sourcetest{5.2.5}{117}{125}
\problemtext 在 $p$--$V$ 图中，分别用何种几何对象表征平衡态、准静态过程和循环过程？
\workings 一个平衡态用\emph{一个点}；一段准静态过程用\emph{一条曲线}；闭合准静态循环用\emph{一条闭合曲线}。
\crosscheck 非准静态过程的中间态通常无法赋予唯一的系统压强。

\sourcetest{5.2.6}{117}{125}
\problemtext 题5.2.2图中理想气体沿 $a\to b$ 等容升压、再 $b\to c$ 等压膨胀；判断气体吸热 $Q$ 与内能变化 $\Delta U$ 的符号。
\workings 两段均使 $T\propto pV$ 上升，$\Delta U>0$；第二段气体对外做功 $W>0$、首段 $W=0$，由 $Q=\Delta U+W$ 知 $\boxed{Q>0,\ \Delta U>0}$。
\crosscheck 初末温度不同，图中经过原点的虚线仅为状态标尺，不是等温线。

\sourcetest{5.2.7}{117}{125}
\problemtext 1kg水由30°C升至80°C，定压比热约4.2J/(g K)，同期电流做功 $4.2\times10^5$J；求水从外界通过热传递获得的热量（正负号）。
\workings 水内能增加近似 $\Delta U=1000\cdot4.2\cdot50=2.1\times10^5$J；电功输入 $W_e=4.2\times10^5$J，故传热 $Q=\Delta U-W_e=\boxed{-2.1\times10^5\ {\rm J}}$。
\crosscheck 负号表示水向环境净散热；要区分传热与电能输入。

\sourcetest{5.2.8}{117}{125}
\problemtext 1mol刚性双原子理想气体在等压升温1K中，内能增加20.78J；求对外做功和吸热。
\workings 由理想气体等压功 $W=nR\Delta T=\boxed{8.31\ {\rm J}}$；第一定律得 $Q=\Delta U+W=20.78+8.31=\boxed{29.09\ {\rm J}}$。
\crosscheck 该20.78J对应五自由度的 $5R/2$ 热容，和题给内能变化一致。

\sourcetest{5.2.9}{117}{125}
\problemtext 3mol单原子理想气体，初态 $p_1=6$atm、$T=500$K，等温膨胀至 $p_2=3$atm；求吸热。
\workings 等温理想气体 $\Delta U=0$、$V_2/V_1=p_1/p_2=2$，故 $\boxed{Q=W=nRT\ln2=3(8.31)(500)\ln2\simeq8.64\times10^3\ {\rm J}}$。
\crosscheck 结果与单原子自由度无关，因理想气体等温内能不变。
